\section{Contribución dos: Dialogue-Aware Search y piKL}

Además del lenguaje controlable, hace falta un planner capaz de convivir con humanos. El self-play puro puede buscar un alto rendimiento, pero puede converger a convenciones que ningún humano ha visto; la imitation pura resulta fácil de entender, pero solo copia los errores del humano promedio y puede ser manipulada con un único mensaje inductor. piKL coloca estos dos extremos dentro de un mismo objetivo de optimización.

\subsection{De la lista de contribuciones al problema de planificación}

La slide 31 resalta la segunda contribución: el planner debe considerar a la vez el dialogue y el human behavior. No puede tratar los mensajes como texto irrelevante, porque un compromiso creíble cambia las acciones de los demás; tampoco puede obedecer los mensajes de forma incondicional, porque un oponente puede proponer un plan desfavorable para Cicero.

\lecturefigure{slide-31-main-contribution-planning.jpg}{Main contribution: planning orientado al diálogo y al comportamiento humano}{Diapositiva V031 recuperada de la grabación oficial; Stanford CS25 Lecture 14}

\subsection{Fallas complementarias de search e imitation}

En el lado izquierdo de la slide 32, la pure search / RL persigue el return, pero no está sujeta a las convenciones humanas; en el lado derecho, la imitation se ciñe a los datos, pero hereda acciones débiles y aprende las correlaciones del lenguaje como correlaciones superficiales. En el centro, piKL funciona como un «pepinillo» que une ambos extremos: toma la human policy como anchor y luego permite que la search se aparte de ella guiada por el valor.

\lecturefigure{slide-32-dialogue-aware-search-rl.jpg}{Dialogue-aware search: buscar una estrategia utilizable entre la optimización pura y la imitación pura}{Diapositiva V032 recuperada de la grabación oficial; Stanford CS25 Lecture 14}

\begin{warningbox}{Cómo falla cada uno de los dos extremos}
Una policy de self-play pura puede adoptar convention que solo entienden sus propias copias, lo que impide coordinarse con humanos; una policy de imitation pura, en cambio, puede reproducir fielmente el comportamiento subóptimo de los humanos y reaccionar en exceso a mensajes como «cédeme tu territorio clave». Un agent capaz de comunicarse debe resultar predecible para los humanos y, al mismo tiempo, conservar la capacidad de rechazar malas sugerencias.
\end{warningbox}

\subsection{piKL: un equilibrio explícito entre el valor y el prior humano}

Sea $\pi_H(a\mid s,h,d)$ la anchor policy aprendida de datos humanos y $Q^{\pi}(s,a)$ el valor de la acción bajo search; el objetivo intuitivo de piKL para un solo estado puede escribirse como
\begin{equation}
\pi^*
=\arg\max_{\pi}
\left[
\mathbb{E}_{a\sim\pi}Q^{\pi}(s,a)
-\lambda D_{\mathrm{KL}}\!\left(\pi(\cdot\mid s)\,\|\,\pi_H(\cdot\mid s)\right)
\right].
\end{equation}
$\lambda$ controla el costo de apartarse del comportamiento humano. Si $\lambda$ es demasiado pequeño, la estrategia puede ser fuerte pero extraña; si es demasiado grande, el sistema se acerca a la imitation original y difícilmente corrige las acciones subóptimas presentes en los datos. No trata la human-likeness como objetivo final, sino que usa las convention humanas como prior de planificación y como modelo de lo que esperan los demás.

\lecturefigure{slide-33-pikl.jpg}{piKL: KL-regularized search anclada en la human imitation policy}{Diapositiva V033 recuperada de la grabación oficial; Jacob et al., ICML 2022}

\begin{knowledgebox}{Lectura de la figura: el lugar de tres tipos de modelo dentro del mismo ciclo}
A la izquierda de la figura, la human policy proporciona el anchor; en el centro, la search / RL calcula la mejora de valor; a la derecha, la retroalimentación del entorno actualiza el estado. La regularización KL no es un simple post-processing, sino que limita de forma continua la deriva de la distribución en cada actualización de la estrategia. Los resultados deben evaluar a la vez la exploitability / return y la human predictability; no basta con reportar solo uno de los dos lados.
\end{knowledgebox}

\teachervoice{El ponente señala que una KL penalty de este tipo no es novedosa; lo que distingue a Cicero es que usa la anchor policy al mismo tiempo como modelo de «qué haría un humano» y de «qué cree un humano que harán los demás». Aquí aparece un matiz de common knowledge / theory-of-mind: el valor de una acción depende de que los demás la consideren una convention razonable.}

\subsection{Cómo cambia realmente el diálogo la estrategia}

La slide 34 muestra en secuencia los cambios de estrategia tras tres tipos de diálogo. En el primero, el mensaje indica que el aliado se retirará de cierta dirección, por lo que Cicero regresa a defender; el segundo expresa desconfianza, lo que lo lleva a reducir las acciones que dependen de esa cooperación; el tercero propone algo desfavorable para Cicero, y el value model impide que el sistema lo siga ciegamente. Lo esencial no es que «cada mensaje cambie la acción», sino que solo los mensajes creíbles y con consecuencias estratégicas deben modificar la distribución de probabilidad.

\lecturefigure{slide-34-strategic-response-sequence.jpg}{Strategic response sequence: cómo las señales del diálogo cambian, o no, la policy de acción}{Diapositiva V034 recuperada de la grabación oficial; Stanford CS25 Lecture 14}

\begin{importantbox}{Dialogue-aware no equivale a instruction-following}
Los mensajes del oponente son una observation con fines estratégicos, no instrucciones del sistema. El planner debe tratarlos como evidencia para actualizar su belief y luego decidir con su propia value function si responde a ellos. Esta distinción permite que Cicero coopere y que también rechace la manipulation.
\end{importantbox}

\subsection{Modeling Human Policies}

La sección anterior mostró que la policy se desplaza con el diálogo creíble; para lograr esa actualización, el sistema primero debe convertir el texto en predicciones de comportamiento. Esta sección explica cómo el planner consume el lenguaje, y no solo cómo lo genera. La slide 35 alimenta a BART con el historial en lenguaje natural para predecir las orders humanas. Esto ilustra una interfaz que a menudo se pasa por alto: el modelo de lenguaje no solo sirve para generar mensajes, sino también para convertir el diálogo en predicciones probabilísticas del comportamiento de los demás jugadores. Después, la distribución predicha entra en la search, lo que forma una ruta con capacidad de auditoría que va del dialogue a la action.

\lecturefigure{slide-35-model-human-policies.jpg}{Modeling human policies: predecir la distribución de acciones humanas a partir del diálogo y del estado}{Diapositiva V035 recuperada de la grabación oficial; Stanford CS25 Lecture 14}

\teachervoice{El docente muestra con varios ejemplos que Cicero no cambia su plan solo porque el otro «lo dijo». La policy se mueve de forma notable únicamente cuando el mensaje, la situación, el historial y los intereses del otro jugador respaldan en conjunto cierto comportamiento; el value estimate, por su parte, se encarga de bloquear las incitaciones que podrían aprovecharse fácilmente del imitation model.}

\subsection{Resumen de la sección}

Mediante la regularización KL, piKL combina la fuerza estratégica de la search con la capacidad de coordinación que dan las convenciones humanas. La esencia de la planificación con conciencia del diálogo tampoco es obedecer al lenguaje, sino convertirlo en una distribución de las acciones de los demás y, después, actualizar la propia estrategia bajo restricciones de valor.
